% 张三慧第四版V2第29章《固体中的电子》原题29.1--29.14独立LaTeX解答
% 主书PDF452--453，S PDF386--389；仅候选，科研/学术专家复核待办。
\section*{第29章\quad 固体中的电子：独立解析}
\paragraph{习题 29.1}
三维自旋二分之一自由电子气零温费米能量
\(E_F=\hbar^2(3\pi^2n)^{2/3}/(2m_e)\)。
每个金原子一个自由电子，
\(n=\rho N_A/M\)，故
\[
n=\frac{19.3\times10^3}{0.197}
(6.02214076\times10^{23})
\simeq5.90\times10^{28}\,\mathrm{m^{-3}}.
\]
于是
\[
\boxed{E_F\simeq5.53\,\mathrm{eV}},\
\boxed{v_F=\sqrt{2E_F/m_e}\simeq1.39\times10^6\,\mathrm{m/s}},
\]
\[
\boxed{T_F=E_F/k_B\simeq6.41\times10^4\,\mathrm K},\
\boxed{\lambda_F=h/\sqrt{2m_eE_F}\simeq0.522\,\mathrm{nm}}.
\]
这属于各向同性抛物能带的简单金属电子气模型，
金的真实能带、有效质量及多能带效应未计。

\paragraph{习题 29.2}
含两个自旋态的三维自由电子能态密度
\(g(E)=C E^{1/2}\)（每单位体积，每能量单位），
在零温所有 \(0<E<E_F\) 状态都被占据：
\[
n=\int_0^{E_F}g(E)\,dE
=C\frac23E_F^{3/2},
\]
\[
u=\int_0^{E_F}E\,g(E)\,dE
=C\frac25E_F^{5/2}.
\]
两式相除：
\[
\boxed{u=\frac35nE_F},\qquad
\boxed{\bar E=u/n=\frac35E_F}.
\]
不能把零温等同于电子完全没有动能；
Pauli填充带来非零的费米动能。

\paragraph{习题 29.3}
对零温自旋填满的费米球，
占据速度分布 \(dN\propto v^2dv\),
\(0\le v\le v_F\)。
所以
\[
\bar v=\frac{\int_0^{v_F}v(v^2dv)}
{\int_0^{v_F}v^2dv}
=\boxed{\frac34v_F},
\]
\[
\langle v^2\rangle=
\frac{\int_0^{v_F}v^2(v^2dv)}
{\int_0^{v_F}v^2dv}
=\frac35v_F^2
\quad\Longrightarrow\quad
\boxed{v_{\rm rms}=\sqrt{\frac35}\,v_F}.
\]
这是零温费米速率\emph{分布的均方根}，
不是单个粒子在Fermi表面的速率。

\paragraph{习题 29.4}
非相对论中子气以质量密度
\(\rho=5\times10^{16}\,\mathrm{kg/m^3}\) 对应
\(n=\rho/m_n\)，中子质量
\(m_n=1.675\times10^{-27}\,\mathrm{kg}\)。
\[
E_{F,n}=\frac{\hbar^2}{2m_n}(3\pi^2n)^{2/3},
\quad
v_{F,n}=\sqrt{2E_{F,n}/m_n}.
\]
代入：
\[
\boxed{E_{F,n}\simeq1.91\times10^7\,\mathrm{eV}
=19.1\,\mathrm{MeV}},\qquad
\boxed{v_{F,n}\simeq6.04\times10^7\,\mathrm{m/s}
\simeq0.20c}.
\]
因为 \(v_F/c\simeq0.20\)，非相对论近似有
百分之几修正，但更重要的是\emph{强核相互作用、
引力场、物态方程}的复杂性；
不能将此自由气模型当作中子星精密物态模型。

\paragraph{习题 29.5}
低温费米气电子摩尔热容
\(C_{V,e}=\pi^2Nk_B(k_BT)/(2E_F)\)。
与经典三维单原子气热容
\(C_{V,\rm classical}=(3/2)Nk_B\) 之比为
\[
\frac{C_{V,e}}{C_{V,\rm classical}}
=\frac{\pi^2k_BT}{3E_F}.
\]
设此比值为0.10且 \(E_F=5\,\mathrm{eV}\)：
\[
\boxed{T=\frac{.3 E_F}{\pi^2k_B}
\simeq1.76\times10^3\,\mathrm K}.
\]
计算时 \(k_BT/E_F\simeq0.0304\)，
仍在Sommerfeld低温相对费米温度的适用区间；
若材料熔化或发生相变则模型可能失效。

\paragraph{习题 29.6}
写钾的摩尔热容为
\(C_m=\gamma T+\beta T^3\)，其中
\(\gamma=2.08\times10^{-3}\),
\(\beta=2.57\times10^{-3}\)
以相应SI摩尔热容单位计。
Sommerfeld系数
\(\gamma=\pi^2Rk_B/(2E_F)\)，故
\[
\boxed{E_F=\frac{\pi^2Rk_B}{2\gamma}
\simeq1.70\,\mathrm{eV}}.
\]
电子热容与晶格热容贡献相等满足
\(\gamma T=\beta T^3\)（\(T>0\)），所以
\[
\boxed{T_*=\sqrt{\gamma/\beta}
\simeq0.900\,\mathrm K}.
\]
此模型按一摩尔钾一个传导电子假设，
若载流电子数不同须相应调整。

\paragraph{习题 29.7}
银按每原子一个自由电子，
\[
n=\frac{(10.5\times10^3)N_A}
{0.1078682}\simeq5.86\times10^{28}\,\mathrm{m^{-3}}.
\]
Drude电阻率
\(\rho_{\rm el}=m_e/(ne^2\tau)\)，得
\[
\boxed{\tau=\frac{m_e}{ne^2\rho_{\rm el}}
\simeq3.78\times10^{-14}\,\mathrm s}.
\]
(2) 300K Maxwell经典平均速率
\(\bar v_{\rm cl}=\sqrt{8k_BT/(\pi m_e)}
\simeq1.08\times10^5\,\mathrm{m/s}\)，
\[
\boxed{\ell_{\rm cl}=\bar v_{\rm cl}\tau
\simeq4.07\,\mathrm{nm}}.
\]
(3) 费米速率
\(v_F=\hbar(3\pi^2n)^{1/3}/m_e
\simeq1.39\times10^6\,\mathrm{m/s}\)，
\[
\boxed{\ell_F=v_F\tau\simeq52.6\,\mathrm{nm}}.
\]
(4) 每离子体积约 \(1/n\)，
平均间距
\[
\boxed{d\sim n^{-1/3}\simeq0.257\,\mathrm{nm}}.
\]
所以 \(\ell_{\rm cl}/d\simeq16\)，
\(\ell_F/d\simeq205\)。经典和费米两种估算的
不同来自真实金属的传导电子不是Maxwell气体；
在简并气体应优先用费米表面速度。

\paragraph{习题 29.8}
费米--狄拉克占据概率
\(f(E)=[\exp((E-E_F)/(k_BT))+1]^{-1}\)。
\(T=1000\,\mathrm K\)，则
\(k_BT\simeq0.08617\,\mathrm{eV}\)。
\[
\boxed{f(E_F+0.10\,\mathrm{eV})
=\frac1{e^{0.10/(k_BT)}+1}\simeq0.239},
\]
\[
\boxed{f(E_F-0.10\,\mathrm{eV})
=\frac1{e^{-0.10/(k_BT)}+1}\simeq0.761}.
\]
两值之和为1，符合分布关于化学势的对称性。
注意这是\emph{每单粒子态}的平均占据数，
不是实际整个能量区间的电子绝对总数。

\paragraph{习题 29.9}
金刚石带隙
\(E_g=5.5\,\mathrm{eV}\)，室温300K下，
在不计态密度差和费米能级偏移的简化Boltzmann近似：
\[
\boxed{\frac{N_{\rm upper}}{N_{\rm lower}}
\simeq e^{-E_g/(k_BT)}
\simeq4.0\times10^{-93}}.
\]
跨越全带隙的阈值波长
\[
\boxed{\lambda_{\max}=\frac{hc}{E_g}
\simeq225\,\mathrm{nm}},
\]
处于紫外光区。这并不表示晶体对所有低于带隙的光
都无任何其他吸收：还可能有缺陷、激子等结构。

\paragraph{习题 29.10}
忽略补偿，每引入一个五价磷原子贡献一个电子。
原本 \(n_0=10^{16}\,\mathrm{m^{-3}}\)，
最终增加到 \(10^6n_0\)，
掺杂电子密度
\(\Delta n\simeq(10^6-1)n_0\simeq10^{22}\,\mathrm{m^{-3}}\)。
硅原子密度
\[
n_{\rm Si}=\frac{2.33\times10^3}
{28.085\times10^{-3}}N_A
\simeq5.00\times10^{28}\,\mathrm{m^{-3}}.
\]
因此需要替换比例
\[
\boxed{x=\Delta n/n_{\rm Si}\simeq2.00\times10^{-7}}
\]
（约每五百万硅原子换一个磷原子）。
对于1.0g硅，等效掺入磷质量
\[
\boxed{m_{\rm P}
\simeq\frac{1.0\,\mathrm g}{28.085}
(30.974)x\simeq2.21\times10^{-7}\,\mathrm g
=0.221\,\mathrm{\mu g}}.
\]
计算是注入磷原子的质量，而不等于替换硅原子后
样品净增加质量。

\paragraph{习题 29.11}
掺磷后的施主能级在导带底下
\(\Delta E_D=0.045\,\mathrm{eV}\)。
光电离施主电子进入导带的最小光子能量是
\(\Delta E_D\)，不是硅的整段
\(E_g=1.2\,\mathrm{eV}\)：
\[
\boxed{\lambda_{\max}=\frac{hc}{\Delta E_D}
\simeq\frac{1239.842}{0.045}\,\mathrm{nm}
=27.6\,\mathrm{\mu m}}.
\]
该红外阈值模型忽略施主束缚态和声子辅助跃迁。

\paragraph{习题 29.12}
光吸收激发跨越禁带需
\(hc/\lambda\ge E_g\)：
\[
\boxed{\lambda_{\max}(\mathrm{CdS})
=\frac{1239.842\,\mathrm{eV\,nm}}{2.42\,\mathrm{eV}}
\simeq512\,\mathrm{nm}\quad\text{可见光}},
\]
\[
\boxed{\lambda_{\max}(\mathrm{PbS})
=\frac{1239.842}{0.30}\,\mathrm{nm}
\simeq4.13\,\mathrm{\mu m}\quad\text{红外光}}.
\]
这里为直接使用带隙估算的本征光电导阈值，
器件实际光谱响应还取决于吸收系数等。

\paragraph{习题 29.13}
Ga--As--P二极管带隙
\(E_g=1.90\,\mathrm{eV}\)，故
\[
\boxed{\lambda_{\max}\simeq\frac{hc}{E_g}
=\frac{1239.842}{1.90}\,\mathrm{nm}
\simeq653\,\mathrm{nm}}.
\]
处于红色可见光区。
真实发射光谱可能因材料组分、缺陷、
温度和直接/间接禁带而变。

\paragraph{习题 29.14}
140nm光子能量
\[
\boxed{E_\gamma=\frac{hc}{140\,\mathrm{nm}}
\simeq8.86\,\mathrm{eV}}.
\]
由于 \(8.86\,\mathrm{eV}>E_g=7.6\,\mathrm{eV}\)，
光子能够激发价带电子跨禁带至导带，
所以按单光子带间吸收模型
\[
\boxed{\text{KCl对140nm紫外光不透明（可强吸收）}}.
\]
有限厚度、选择定则和表面反射会影响具体透射率，
本题只问带隙阈值判定。
